\SourceLeaf{402}%
\par
Ist \(t\) die komplexe Zahl für den Schwerpunkt, so gilt \(3t=\alpha+\beta+\gamma\) und damit
\[
\begin{aligned}
3t-2u&=\alpha+(\beta-\gamma)\operatorname{cotg}A.i\\
&=\beta+(\gamma-\alpha)\operatorname{cotg}B.i\\
&=\gamma+(\alpha-\beta)\operatorname{cotg}C.i
\end{aligned}
\]
\par
Dieser Ausdruck \(3t-2u\) ist die komplexe Zahl für den Punkt, in dem sich die drei Senkrechten von den Eckpunkten auf die gegenüberliegenden Seiten schneiden.

Daraus folgt durch Subtraktion
\[
\begin{aligned}
0&=\alpha-\beta+\{\alpha\operatorname{cotg}B+\beta\operatorname{cotg}A-\gamma(\operatorname{cotg}A+\operatorname{cotg}B)\}.i\\
0&=\beta-\gamma+\{\beta\operatorname{cotg}C+\gamma\operatorname{cotg}B-\alpha(\operatorname{cotg}B+\operatorname{cotg}C)\}.i\\
0&=\gamma-\alpha+\{\gamma\operatorname{cotg}A+\alpha\operatorname{cotg}C-\beta(\operatorname{cotg}C+\operatorname{cotg}A)\}.i
\end{aligned}
\]
\par\medskip\Needspace{5\baselineskip}
\pdfbookmark[2]{Zweite handschriftliche Bemerkung.}{carnot-handschriftlich-ii}
\begin{center}[\textit{Zweite handschriftliche Bemerkung.}]\end{center}
\hspace*{\parindent}Es seien \(a,b,c,d\) vier Punkte auf dem Umfang eines Kreises mit dem Radius \(1\) und zugleich die komplexen Zahlen, die diesen Punkten entsprechen [wobei die dem Mittelpunkt entsprechende komplexe Zahl gleich \(0\) gesetzt wird]. Ferner seien \(p,q,r\) die Schnittpunkte der Geraden \(\tfrac{ab}{cd},\tfrac{ac}{bd},\tfrac{ad}{bc}\) und \(p^*\) der Mittelpunkt der Kreissehne, an deren Endpunkten sich die Tangenten in \(p\) schneiden. Ebenso seien \(q^*,r^*\) definiert. Wenn die Buchstaben mit Strich jeweils die zugehörigen konjugierten Zahlen bezeichnen, gilt
\[
\begin{aligned}
p&=\frac{abc+abd-acd-bcd}{ab-cd}=a-\frac{b(a-c)(a-d)}{ab-cd}\\
q&=\frac{abc+acd-abd-bcd}{ac-bd}=a-\frac{c(a-b)(a-d)}{ac-bd}\\
r&=\frac{acd+abd-abc-bcd}{ad-bc}=a-\frac{d(a-b)(a-c)}{ad-bc}
\end{aligned}
\]
\[
p-q=(a-d)(b-c)\frac{abd+acd-abc-bcd}{(ab-cd)(ac-bd)}=\frac{(a-d)(b-c)(ad-bc)r}{(ab-cd)(ac-bd}
\]
\[
p^*=\frac1{p'}=\frac{ab-cd}{a+b-c-d}=a-\frac{(a-c)(a-d)}{a+b-c-d}
\]
\begingroup\footnotesize
\[
p^*-q=\frac{(a-d)(b-c)(ad-bc)}{(ac-bd)(a+b-c-d)}=p^*.\frac{(a-d)(b-c)(ad-bc)}{(ab-cd)(ac-bd)}=p^*\frac{p-q}{r}
\]
\endgroup
oder \(p^*(q+r-p)=qr\), ebenso \(q^*(p+r-q)=pr\) und \(r^*(p+q-r)=pq\)
